% ============================================================
% RESUMEN
% ============================================================
\chapter*{Resumen}
\addcontentsline{toc}{chapter}{Resumen}
\markboth{RESUMEN}{RESUMEN}

La técnica fotoacústica pulsada permite medir coeficientes de absorción
óptica varios órdenes de magnitud menores que los accesibles por
espectrofotometría de transmisión, lo que la hace apropiada para
líquidos que son esencialmente transparentes a la longitud de onda de
excitación. Su costo es operativo: cada punto experimental exige
coordinar el disparo del láser, la adquisición del osciloscopio y el
registro de la temperatura de la muestra, una coordinación que
realizada a mano resulta lenta y difícil de reproducir entre sesiones.

Este trabajo desarrolla el sistema de automatización de un arreglo
fotoacústico del Laboratorio de Fotofísica y Películas Delgadas del
ICAT. El sistema integra tres instrumentos bajo una sola aplicación de
escritorio escrita en Python con interfaz gráfica PySide6: un láser
Nd:YAG pulsado EKSPLA NL303HT-10-SH, controlado a través de la
biblioteca de enlace dinámico del fabricante; un osciloscopio digital
Tektronix de la serie TDS5000B, comunicado por VXI-11 sobre Ethernet; y
un módulo de medición de temperatura desarrollado íntegramente en este
trabajo, compuesto por un microcontrolador ESP32 WROOM-32 y cuatro
sensores DS18B20 sobre buses 1-Wire independientes.

La aplicación ejecuta secuencias de medición sin intervención del
operador en dos modos: adquisición durante un intervalo de tiempo fijo,
y adquisición gobernada por la temperatura de la muestra, en la que el
osciloscopio opera como integrador entre los instantes en que el
líquido cruza umbrales sucesivos durante su enfriamiento. Cada captura
se guarda con las muestras crudas del instrumento y con los metadatos
necesarios para reconstruirla en unidades físicas, incluidos los
parámetros del láser, las escalas vigentes del osciloscopio y una
marca explícita de las condiciones anómalas detectadas durante la
adquisición.

La validación se organizó en dos planos. La lógica del programa se
ejercitó sin instrumentos mediante simulación de perfiles térmicos,
recorrido completo de la cadena de medición y almacenamiento, y
sesiones sintéticas construidas para provocar cada una de las
condiciones que la herramienta de verificación debe detectar; esta
última se sometió además a una prueba de mutación que confirmó que
ninguna de sus comprobaciones puede desactivarse sin que la prueba
correspondiente falle. El comportamiento del sistema integrado se
evaluó sobre sesiones reales de laboratorio.

El criterio con el que se juzga el resultado no es la ausencia de
fallas, inevitable en un sistema que integra tres instrumentos
independientes, sino la capacidad del sistema para dejar constancia de
ellas en lugar de aparentar normalidad. Se documentan cinco anomalías
observadas durante el desarrollo, clasificadas según el nivel al que
fue posible detectarlas —durante la ejecución, en la verificación
posterior de los datos, o en la auditoría del código—, incluida una que
tuvo como origen la propia herramienta de verificación. Cuatro de las
cinco no eran observables mientras ocurrían, y todas resultaron
reconstruibles a partir de la información conservada en los archivos de
sesión. Esa reconstruibilidad es la propiedad que el sistema aporta al
arreglo experimental.
